\begin{exercice}\label{exo4-1} (\minsyndical)

Soit une urne qui contient 3 jetons noirs et 2 jetons blancs. Un joueur tire simultan\'ement 2 jetons de l'urne. S'il tire 2 jetons blancs, il gagne 10 euros, s'il tire 2 jetons noirs, il gagne 2 euros, sinon il perd 5 euros (quand donc il tire un jeton noir et un blanc).

\begin{enumerate}

\item D\'ecrire explicitement l'univers $\Omega$ associ\'e \`a l'\'epreuve.

\item Soit $X : \Omega \rightarrow \mathbb{R}$ la variable al\'eatoire associ\'ee au gain, qui peut donc prendre pour valeurs $-5$ , $2$ et $10$. Donner la loi de probabilit\'e associ\'ee \`a $X$.

\item Calculer l'esp\'erance math\'ematique, la variance et l'\'ecart-type de $X$.

\item Le jeu est-t'il \'equitable ? Proposer une r\`egle de r\'epartition du gain qui le rende \'equitable.\\

\end{enumerate}

%\corrref{TD1_1}

\end{exercice}
